\originalchapter{带电粒子与氢原子的量子对称性}
本章为补足已选进阶内容的独立推导. 8.06 Spring 2018没有公开第5章带电粒子PDF, 所以这里没有冒称翻译该缺失原件. 公式按前文算符、概率流和角动量工具组织, 带电粒子例子与因子分解方法也为考核篇提供所需支持.
\originalsection{电磁势、机械动量与规范}
在电场中粒子可受力、在磁场中方向可偏转. 用标量势$\phi$和矢势$\boldsymbol A$表示$\boldsymbol E=-\nabla\phi-\partial_t\boldsymbol A$,$\boldsymbol B=\nabla\times\boldsymbol A$. 对质量$m$、电荷$q$的无自旋粒子, 最小耦合Hamiltonian为$H=(\boldsymbol p-q\boldsymbol A)^2/(2m)+q\phi$, $\boldsymbol p=-\ii\hbar\nabla$为正则动量, $\boldsymbol\pi=\boldsymbol p-q\boldsymbol A$为机械动量. 平方有算符顺序, 要按$(p_i-qA_i)(p_i-qA_i)$作用, 不能先把$p_i,A_i$交换. $q\boldsymbol A$量纲为动量.

同一电磁场可以用不同势: $\boldsymbol A' =\boldsymbol A+\nabla\chi$, $\phi'=\phi-\partial_t\chi$. 波函数同时变$\psi'=\e^{\ii q\chi/\hbar}\psi$, 直接求导得$\boldsymbol\pi'\psi'=\e^{\ii q\chi/\hbar}\boldsymbol\pi\psi$, 时间导数中额外$-q\dot\chi$与标势变化抵消, 所以方程协变. 此处$\chi$量纲作用量/电荷. 密度不变, 流为$\boldsymbol J=m^{-1}\operatorname{Re}(\psi^*\boldsymbol\pi\psi)$, 也保持不变. 规范选择改变描述中的局部相位, 不改变电磁场或计数统计.

机械动量分量一般不对易: $[\pi_i,\pi_j]=\ii q\hbar\epsilon_{ijk}B_k$, 由$[p_i,A_j]=-\ii\hbar\partial_iA_j$求得. 所以不能在磁场中继续将三个机械动量当成自由粒子相容读数. 用Heisenberg方程对位置求导给速度$\boldsymbol\pi/m$, 可与经典Lorentz力联系; 非均匀场时算符乘积要对称排列, 均值也不总能以经典点轨迹代替.
\originalsection{均匀磁场的Landau能级}
取$B\hat z$, 无电势, 选$\boldsymbol A=(0,Bx,0)$. $p_y,p_z$与$H$对易, 对其广义本征值$\hbar k_y,p_z$将波函数写成$\e^{\ii k_yy+\ii p_zz/\hbar}u(x)$. 剩下
\[
H_x=\frac{p_x^2}{2m}+\frac{(qB)^2}{2m}(x-x_c)^2,\quad
x_c=\frac{\hbar k_y}{qB},\quad\omega_c=\frac{|qB|}{m},
\]
故$E_{n,p_z}=\hbar\omega_c(n+1/2)+p_z^2/(2m)$,$n\geq0$. 新能量不依赖中心$x_c$, 同一能级有很多轨道中心制备; 这不是漏掉变量而是磁平移对称导致简并. 波函数宽度尺度$\ell_B=\sqrt{\hbar/|qB|}$, 能级间隔$\hbar\omega_c$对应经典回旋频率, 但最低能仍有零点能量.

若矩形面积$L_xL_y$, $y$周期使$k_y=2\pi j/L_y$, 允许中心在样品内, 每个Landau能级约有$|qB|L_xL_y/(2\pi\hbar)$个态, 每面积态数$|qB|/(2\pi\hbar)$. 边界附近能量可变化, 该数忽略边缘有限尺寸修正. 若为电子还需自旋项$+\mu_BB\sigma_z$（$g\simeq2$, 不计辐射异常磁矩）, 与轨道能量相加; 无自旋结果不是完整电子磁谱.
\originalsection{无局部磁场区域的Aharonov--Bohm相位}
细长磁通管外部$\boldsymbol B=0$, 但绕管闭合路径$\oint\boldsymbol A\cdot\dd\boldsymbol l=\Phi$. 两条绕行路径的振幅相对相位为$q\Phi/\hbar$, 可改变干涉计数. 在单连通的小区域可用规范移除$\boldsymbol A$, 围住被排除磁通管的全区域却未必可用单值规范同时移除; 边界拓扑参与条件, 不代表粒子在局部无场区域受经典磁力.

半径$b$环上的粒子是可算模型. 取$A_\varphi=\Phi/(2\pi b)$, 单值边界$\psi(\varphi+2\pi)=\psi(\varphi)$给$\psi_n=\e^{\ii n\varphi}/\sqrt{2\pi}$,$n\in\mathbb Z$（按角度归一）. 能量
\[
E_n(\Phi)=\frac{\hbar^2}{2mb^2}\left(n-\frac{q\Phi}{h}\right)^2.
\]
磁通改变机械角动量, 能量谱作为集合的周期为$h/|q|$, 因磁通加一个量子只重标$n$. 固定$n$曲线不是周期函数, 需区分单支与全谱. 若把矢势变零, 波函数边界变扭转$\psi'(2\pi)=\e^{-\ii q\Phi/\hbar}\psi'(0)$; 同时偷偷保留原周期边界会改变物理问题.
\originalsection{径向因子分解与氢谱}
第11章由边界和级数求氢谱; 此处给另一条不依赖非终止级数渐近的代数路线. 本节取固定核$m=m_e$,$a_Z=\hbar^2/(mg)$; 若计核运动则同时将$m$改为约化质量$\mu$, 不混用两套长度和质量. 对固定$\ell$的半轴Coulomb径向算符, 取实超势（只是因子分解的函数名, 不是另一个物理势）$W_\ell=-(\ell+1)/r+1/[(\ell+1)a_Z]$, $A_\ell=\mathrm d/\mathrm dr+W_\ell$, $A_\ell^\dagger=-\mathrm d/\mathrm dr+W_\ell$, 使用正规径向边界使伴随积分项消失. 展开得
\[
H_\ell=\frac{\hbar^2}{2m}A_\ell^\dagger A_\ell-\frac{\hbar^2}{2ma_Z^2(\ell+1)^2}.
\]
正性给基态下界, $A_\ell u=0$解$u\propto r^{\ell+1}\e^{-r/[(\ell+1)a_Z]}$确实正规且取到下界. 伙伴$A_\ell A_\ell^\dagger$将离心项变为$(\ell+1)(\ell+2)/r^2$, 对应$H_{\ell+1}$加已知常数. 用$H^{(1)}A^\dagger=A^\dagger H^{(2)}$逐次从更大$\ell$的零态提升, 得$E_{n_r,\ell}=-\hbar^2/[2ma_Z^2(n_r+\ell+1)^2]$, 与第11章一致. 因子分解不是每个势都具有的结构; 一旦零态不正规或伙伴不属于同一边界域, 形式阶梯就不能直接报告物理谱.
\originalsection{氢的额外简并与Runge--Lenz量}
一般中心势只有$m$的$2\ell+1$简并, 氢同$n$还包含不同$\ell$, 不由普通空间旋转单独解释. 取固定核$H=p^2/(2m)-g/r$, 定义无量纲$\boldsymbol R=(\boldsymbol p\times\boldsymbol L-\ii\hbar\boldsymbol p)/(mg)-\boldsymbol r/r$. 等价用$(\boldsymbol p\times\boldsymbol L-\boldsymbol L\times\boldsymbol p)/2$使Hermitian性明确. 在合适共同域的Coulomb算符计算给$[H,R_i]=0$, $[L_i,R_j]=\ii\hbar\epsilon_{ijk}R_k$, $[R_i,R_j]=-2\ii\hbar H\epsilon_{ijk}L_k/(mg^2)$以及
\[
\boldsymbol R\cdot\boldsymbol L=0,\qquad R^2=1+\frac{2H}{mg^2}(L^2+\hbar^2).
\]
这是量子修正后的算符恒等式, 可用$[p_i,r_j]$及径向导数逐项检验; $\hbar^2$项不在经典公式中. 垂直性的一段具体计算见2012期末第3题. 这些多项对易恒等式作为本节引用的Coulomb代数工具, 后面的谱推理从它们明确展开, 不冒称已经逐项证明所有长代数式.

固定$E<0$能量空间, 令$\boldsymbol A=\sqrt{-mg^2/(2E)}\boldsymbol R$, $\boldsymbol J_\pm=(\boldsymbol L\pm\boldsymbol A)/2$. 代上式给两套互相对易的角动量代数, $J_+^2=J_-^2=[-mg^2/(2E)-\hbar^2]/4$. 两套量子数相同为$j$, 所以$-mg^2/(2E)=\hbar^2(2j+1)^2$. 记$n=2j+1$, 得氢能量与忽略自旋时简并$n^2$. 不同$\ell=0,\ldots,n-1$维数之和$\sum(2\ell+1)=n^2$, 正是这套隐藏对称性的同一空间. 精细结构或外场改变$H$, 不再自动保持$\boldsymbol R$守恒, 因而能够打破这些简并.
